\section{Number of Distinct Substrings in a String}
\label{sec:trie-distinct-substrings}

\problemstatement{Given a string $s$ of length $n$, count the number of
\textbf{distinct} (non-empty) substrings of $s$.}

\begin{patternbox}
\emph{``Count distinct substrings''} reframes naturally once you notice
that every substring of $s$ is, by definition, a \textbf{prefix of some
suffix} of $s$. Inserting \emph{every suffix} of $s$ into a trie
therefore inserts every substring of $s$ as some path from the root --
and because a trie physically \textbf{shares} identical prefixes as the
same nodes, the number of \textbf{distinct} substrings is exactly the
number of \textbf{distinct nodes} the trie ends up containing.
\end{patternbox}

\subsection*{Understanding the problem carefully}

Consider inserting all $n$ suffixes of $s$ (the suffix starting at index
$0$, the suffix starting at index $1$, \dots, the suffix starting at
index $n-1$) into a trie, one character at a time, exactly as in
Section~\ref{sec:trie-implement}. Every prefix of every suffix is some
substring of $s$, and every substring of $s$ is the prefix of at least
one suffix (specifically, the suffix that starts where the substring
starts) -- so this process inserts every substring of $s$, with
duplicates naturally \emph{not} creating new nodes, since inserting a
character that already has a corresponding child simply reuses that
existing node.

\begin{ideabox}[Count New Nodes Created, Not Total Insertions]
For each suffix of $s$ (there are $n$ of them, of lengths $n, n-1, \dots,
1$), walk it through the trie character by character. Whenever a
required child does not yet exist, create it \emph{and increment a
running counter of distinct substrings} (since a brand-new node means
this exact substring -- the path from root to this new node -- has never
been encountered before). When a required child already exists, simply
move into it without incrementing the counter (this exact substring was
already counted during an earlier suffix's insertion).
\end{ideabox}

\subsection*{Why counting newly created nodes exactly counts distinct substrings}

Each trie node corresponds to exactly one specific substring -- the
concatenation of characters along the path from the root to that node --
and two different nodes can never correspond to the same substring (a
trie's structure forbids two different paths spelling out the same
string, since shared prefixes are always merged into the same nodes).
Conversely, every substring of $s$ does correspond to \emph{some} node,
because it is a prefix of at least one inserted suffix. So there is a
perfect one-to-one correspondence between distinct substrings of $s$ and
nodes of the trie (excluding the root, which represents the empty
string) -- counting nodes created during insertion is therefore counting
distinct substrings directly, with no separate deduplication step needed.

\subsection*{Algorithm}

\begin{enumerate}
  \item Initialize an empty trie and $\textit{count} \gets 0$.
  \item For each starting index $i$ from $0$ to $n-1$ (each suffix
        $s[i..n-1]$):
  \begin{enumerate}
    \item Walk the suffix's characters from the trie's root. For each
          character: if the required child does not exist, create it
          and increment $\textit{count}$; move into the child.
  \end{enumerate}
  \item Return \textit{count}.
\end{enumerate}

\subsection*{Dry run}

\dryrun{Take $s = \texttt{"ab"}$. Suffixes: \texttt{"ab"},
\texttt{"b"}.}

\begin{itemize}
  \item Insert \texttt{"ab"}: \texttt{a} does not exist at the root --
        create it, $\textit{count}=1$. Move into \texttt{a}.
        \texttt{b} does not exist there -- create it, $\textit{count}=2$.
  \item Insert \texttt{"b"}: \texttt{b} does not exist at the root
        (this is a different node than the \texttt{b} created above,
        which was a child of \texttt{a}, not of the root) -- create it,
        $\textit{count}=3$.
\end{itemize}

Total distinct substrings: $3$ -- namely \texttt{"a"}, \texttt{"ab"},
\texttt{"b"}, matching the direct enumeration of all substrings of
\texttt{"ab"}.

\subsection*{C++ implementation}

\begin{lstlisting}
#include <bits/stdc++.h>
using namespace std;

struct TrieNode {
    TrieNode* children[26] = {nullptr};
};

int countDistinctSubstrings(string s) {
    TrieNode* root = new TrieNode();
    int n = s.size();
    int count = 0;

    for (int i = 0; i < n; i++) {
        TrieNode* node = root;
        for (int j = i; j < n; j++) {
            int idx = s[j] - 'a';
            if (!node->children[idx]) {
                node->children[idx] = new TrieNode();
                count++;
            }
            node = node->children[idx];
        }
    }
    return count;
}

int main() {
    cout << countDistinctSubstrings("ab") << endl;   // 3
    cout << countDistinctSubstrings("aaa") << endl;  // 3 (a, aa, aaa)
    return 0;
}
\end{lstlisting}

\begin{complexitybox}
\textbf{Time Complexity.} Inserting all $n$ suffixes, of total length
$n + (n-1) + \cdots + 1$, gives
\[
  O(n^2)
\]
character insertions in the worst case, each $O(1)$ amortized trie work.
(A suffix automaton or suffix array based approach can solve the same
problem in $O(n \log n)$ or $O(n)$, but the trie-of-all-suffixes
approach shown here is the direct, easy-to-derive $O(n^2)$ solution,
appropriate for moderate string lengths.)
\textbf{Space Complexity.} $O(n^2)$ in the worst case (one trie node per
distinct substring, and a string can have up to $O(n^2)$ distinct
substrings).
\end{complexitybox}

\begin{takeawaybox}
``Count distinct substrings'' is solved by recognizing that a trie's
node-sharing behavior \emph{is} a deduplication mechanism: insert
everything that could possibly create a duplicate (here, every suffix,
since every substring is a prefix of some suffix), and count only the
genuinely new nodes created, never the total insertion attempts.
\end{takeawaybox}
